\chapter{第七章：RAG 检索增强生成}

到了这一章，你已经具备了两块关键积木：一是 LLM 的工作原理，二是 embedding 与向量检索。把这两者拼起来，就是 RAG（Retrieval-Augmented Generation，检索增强生成）。在 2026 年，几乎所有``让模型接入企业私有知识''的系统，都离不开 RAG。面试时只要岗位描述里出现下面这些词中的任意两个，RAG 几乎必考：


\begin{itemize}[leftmargin=2em]
\item knowledge base
\item enterprise search
\item chatbot
\item internal docs
\item grounding
\item hallucination reduction
\end{itemize}

但真正能落地的人并不多。因为大多数人停留在``检索几个 chunk 塞进 prompt''这一步，而没有意识到：\textbf{RAG 不是一个 API，而是一条数据处理、检索、排序、生成、评估的完整流水线}。


\bigskip\hrule\bigskip

\section{7.1 为什么需要 RAG：LLM 的三个天然限制}

\subsection{7.1.1 知识截止（Knowledge Cutoff）}

无论模型多强，它的预训练数据总有时间边界。你问它``本公司上周发布的内部制度更新''，模型天然不知道。RAG 的第一价值，就是给模型补充\textbf{新知识}。


\subsection{7.1.2 幻觉（Hallucination）}

模型在缺少证据时，仍然倾向于给出`\texttt{看起来合理''的答案。这对写作场景还勉强能忍，对客服、法务、医疗、运维则可能直接出事故。RAG 的第二价值，是让模型回答时}`有证可依''。


\subsection{7.1.3 无法直接访问私有数据}

你公司的 Jira、Confluence、PDF 规范、代码仓库、数据库、工单系统，都不在通用模型的预训练集里。RAG 的第三价值，是让模型在不重新训练大模型的前提下访问私域信息。


简化理解：


\begin{codeblock}
没有 RAG：模型 = 聪明但靠记忆作答
有 RAG：模型 = 聪明 + 会查资料再作答
\end{codeblock}

\bigskip\hrule\bigskip

\section{7.2 RAG 架构的演进：从 Naive 到 Modular}

\subsection{7.2.1 Naive RAG}

最基础的 RAG 流程如下：


\begin{codeblock}
用户问题
   │
   ▼
Embedding 查询
   │
   ▼
检索 TopK 文档片段
   │
   ▼
拼接到 Prompt
   │
   ▼
LLM 生成答案
\end{codeblock}

这就是``retrieve -> stuff -> generate''。它能用，但很容易遇到三类问题：


\begin{enumerate}[leftmargin=2em]
\item 检索到不相关片段
\item 片段相关但排序不佳
\item 文档里明明有答案，模型还是答错
\end{enumerate}

\subsection{7.2.2 Advanced RAG}

进阶版 RAG 会在检索前后加入更多组件：


\begin{itemize}[leftmargin=2em]
\item Query Transformation（查询改写）
\item Hybrid Search（稠密 + 稀疏混合）
\item Reranking（重排序）
\item Iterative Retrieval（迭代检索）
\item Context Compression（上下文压缩）
\end{itemize}

一个典型的高级流水线：


\begin{codeblock}
用户问题
   │
   ▼
[Query Rewrite]
   │
   ├─ 原始查询
   ├─ 改写查询
   └─ 多查询扩展
   ▼
[Dense + Sparse Retrieval]
   │
   ▼
[Reranker]
   │
   ▼
[Prompt Builder]
   │
   ▼
LLM Answer
\end{codeblock}

\subsection{7.2.3 Modular RAG}

Modular RAG（模块化 RAG）强调：检索、排序、压缩、回答、评估都应可替换。它非常适合 Agent 平台，因为不同业务线可以共享骨架、替换组件。


例如：


\begin{itemize}[leftmargin=2em]
\item 法务知识库：更重视精确召回与引用
\item 代码问答：更重视结构切块和仓库路径信息
\item 客服 FAQ：更重视延迟和低成本
\end{itemize}

模块化设计的核心思想不是``用更多框架''，而是：\textbf{每个组件都可单独评估和升级}。


\bigskip\hrule\bigskip

\section{7.3 完整 RAG 流水线 ASCII 图}

\begin{codeblock}
                    ┌──────────────────────┐
                    │  Documents / Data    │
                    │ PDF / HTML / MD /代码 │
                    └──────────┬───────────┘
                               │
                               ▼
                    ┌──────────────────────┐
                    │ Loading & Cleaning   │
                    └──────────┬───────────┘
                               │
                               ▼
                    ┌──────────────────────┐
                    │ Chunking             │
                    └──────────┬───────────┘
                               │
                               ▼
                    ┌──────────────────────┐
                    │ Embedding            │
                    └──────────┬───────────┘
                               │
                               ▼
                    ┌──────────────────────┐
                    │ Vector / BM25 Index  │
                    └──────────┬───────────┘
                               │
      用户问题                 │
         │                     │
         ▼                     │
┌──────────────────────┐       │
│ Query Transform      │───────┘
└──────────┬───────────┘
           │
           ▼
┌──────────────────────┐
│ Dense / Sparse Search│
└──────────┬───────────┘
           │
           ▼
┌──────────────────────┐
│ Rerank / Compress    │
└──────────┬───────────┘
           │
           ▼
┌──────────────────────┐
│ Prompt Builder       │
└──────────┬───────────┘
           │
           ▼
┌──────────────────────┐
│ LLM Generate Answer  │
└──────────────────────┘
\end{codeblock}

\bigskip\hrule\bigskip

\section{7.4 文档处理：加载、清洗、切块}

RAG 效果一半来自模型，另一半来自文档处理。很多系统召回差，不是 embedding 模型不够强，而是前处理阶段已经把信息切碎了。


\subsection{7.4.1 Loading：数据源不是只有 PDF}

常见数据源包括：


\begin{itemize}[leftmargin=2em]
\item PDF：制度文件、白皮书、投标资料
\item HTML：帮助中心、官网、Wiki
\item Markdown：技术文档、README、设计文档
\item Code Files：\texttt{.py}、\texttt{.ts}、\texttt{.go} 等源码
\item Structured Data：CSV、数据库记录、工单、聊天记录
\end{itemize}

不同类型文档的处理重点不同：


\begingroup
\scriptsize
\setlength{\tabcolsep}{3pt}
\begin{longtable}{|p{0.31\textwidth}|p{0.31\textwidth}|p{0.31\textwidth}|}
\hline
数据类型 & 关键挑战 & 建议 \\ \hline
PDF & 版面噪声、换行乱 & OCR 后做段落重建 \\ \hline
HTML & 导航栏、脚注、广告 & 只保留正文 DOM \\ \hline
Markdown & 标题层级重要 & 保留 heading 路径 \\ \hline
代码 & 结构边界关键 & 按函数、类、文件树切块 \\ \hline
\end{longtable}
\endgroup

\subsection{7.4.2 Chunking 为什么是灵魂}

切块（Chunking）决定了检索单元。RAG 不是检索整篇文档，而是检索 chunk。chunk 过大，噪声多；过小，信息断裂。


\subsubsection{固定长度切块}

最简单：每 500 tokens 一块，重叠 50 tokens。优点是实现简单；缺点是经常把一个完整定义切成两半。


\subsubsection{递归切块（Recursive Chunking）}

按标题、段落、句子、字符逐层退化，尽量保留自然边界。这是技术文档里非常常见的实用方案。


\subsubsection{语义切块（Semantic Chunking）}

根据句向量或主题变化来决定边界。效果往往更自然，但计算成本更高。


\subsubsection{Agentic Chunking}

让模型参与切块，比如``根据问题回答需要的最小语义单元重组文档''。这适合高价值场景，但成本高，不适合海量基础索引。


\subsection{7.4.3 Chunk 大小与 overlap 的取舍}

下面是一组对技术知识库的常见实验结论（示意但贴近真实工程）：


\begingroup
\scriptsize
\setlength{\tabcolsep}{3pt}
\begin{longtable}{|p{0.19\textwidth}|p{0.19\textwidth}|p{0.19\textwidth}|p{0.19\textwidth}|p{0.19\textwidth}|}
\hline
Chunk Size & Overlap & Recall@5 & Answer Faithfulness & 平均输入成本 \\ \hline
128 tokens & 20 & 0.68 & 0.72 & 低 \\ \hline
256 tokens & 30 & 0.79 & 0.81 & 中 \\ \hline
512 tokens & 50 & 0.83 & 0.84 & 中高 \\ \hline
1024 tokens & 80 & 0.84 & 0.76 & 高 \\ \hline
\end{longtable}
\endgroup

可以看到：更大的 chunk 不一定更好。因为虽然召回率略增，但上下文噪声也增加，生成阶段更容易``看花眼''。


经验规则：


\begin{itemize}[leftmargin=2em]
\item FAQ / 短知识：150-300 tokens
\item 技术设计文档：300-600 tokens
\item 法务/制度文件：500-800 tokens，保留章节标题
\item 代码：按函数/类/文件边界，不要强行固定长度
\end{itemize}

\bigskip\hrule\bigskip

\section{7.5 Retrieval 策略：Dense、Sparse、Hybrid 与更多变体}

\subsection{7.5.1 Dense Retrieval}

Dense retrieval（稠密检索）就是 embedding 检索。优点是语义强，缺点是对精确关键词、编号、错误码不够敏感。


适合：


\begin{itemize}[leftmargin=2em]
\item 用户表达多样
\item 同义改写多
\item 中文语义问答
\end{itemize}

\subsection{7.5.2 Sparse Retrieval}

Sparse retrieval（稀疏检索）代表方法是 BM25、TF-IDF。它对关键字、精确术语、变量名、报错码特别敏感。


适合：


\begin{itemize}[leftmargin=2em]
\item 日志检索
\item API 名称、报错码定位
\item 法律条文编号
\end{itemize}

\subsection{7.5.3 Hybrid Retrieval}

Hybrid retrieval 是生产系统里的主流方案：把 dense 和 sparse 的结果合并。原因非常简单：


\begin{itemize}[leftmargin=2em]
\item dense 擅长``意思相近''
\item sparse 擅长``字面精确''
\end{itemize}

一个很常见的融合公式：


\texttt{score = α * dense\_score + (1 - α) * sparse\_score}


实际项目里 \texttt{α} 常在 \texttt{0.5-0.8} 区间调优。


\subsection{7.5.4 Multi-query Retrieval}

用户原问题往往表达不稳定。例如：


\begin{quote}
``为什么我的订单创建接口有时会超时？''\\
\end{quote}

可以自动扩展成：


\begin{itemize}[leftmargin=2em]
\item 订单创建接口超时原因
\item 数据库写入导致订单超时
\item 下游支付服务慢导致订单接口超时
\end{itemize}

多查询检索常能提高召回，但会增加成本和噪声，所以通常需要 rerank。


\subsection{7.5.5 Parent-child Retrieval}

做法是把大文档切成小块建立索引，但返回时附带其``父文档''或更大上下文。它特别适合：


\begin{itemize}[leftmargin=2em]
\item 长技术设计文档
\item 法规条文
\item 多级标题文档
\end{itemize}

你既能用细粒度 chunk 提高召回，又能在回答阶段给模型更完整上下文。


\bigskip\hrule\bigskip

\section{7.6 Reranking：为什么 TopK 之后还要再排一次}

检索系统的一个现实问题是：初步召回的前 20 条里，通常只有 3 到 8 条真正最相关。于是我们引入 reranker（重排序器）。


\subsection{7.6.1 Cross-Encoder Reranker}

Cross-encoder 会把``query + candidate chunk''一起送入模型评分。相比双塔 embedding，它计算更贵，但排序更准。


\subsection{7.6.2 常见 reranker 选择}

\begingroup
\scriptsize
\setlength{\tabcolsep}{3pt}
\begin{longtable}{|p{0.31\textwidth}|p{0.31\textwidth}|p{0.31\textwidth}|}
\hline
Reranker & 特点 & 适用场景 \\ \hline
Cohere Rerank & API 易用、英文强 & 快速上线 \\ \hline
BGE-Reranker & 开源可自部署，中文不错 & 私有化、中文系统 \\ \hline
Jina Reranker & 多语言可选 & 国际化系统 \\ \hline
\end{longtable}
\endgroup

\subsection{7.6.3 一个实战经验}

对 50 万 chunk 的企业文档库，常见配置是：


\begin{enumerate}[leftmargin=2em]
\item dense 检索 30 条
\item sparse 检索 30 条
\item 合并去重后得到 40-50 条
\item rerank 前 10 条
\item 最终送给 LLM 4-6 条
\end{enumerate}

这通常比`\texttt{直接送 Top3''更稳，也比}`直接送 Top20''更省 token。


\bigskip\hrule\bigskip

\section{7.7 RAG 评估：不评估，系统就只能靠感觉调}

RAG 的难点在于，错误可能发生在多个环节：


\begin{itemize}[leftmargin=2em]
\item 文档没加载进来
\item chunk 切坏了
\item 检索没召回
\item rerank 排错了
\item prompt 没约束好
\item 模型看到了正确上下文却答错
\end{itemize}

因此必须分层评估。


\subsection{7.7.1 Retrieval Metrics}

\subsubsection{Recall@K}

正确文档是否出现在前 K 个结果中。最基础、最重要。


\subsubsection{MRR（Mean Reciprocal Rank）}

正确结果排名越靠前，分数越高。适合看排序质量。


\subsubsection{nDCG}

考虑多级相关性，不只是``对/错''，更适合复杂检索评估。


\subsection{7.7.2 Generation Metrics}

\subsubsection{Faithfulness}

答案是否忠于给定上下文，而不是自己瞎编。


\subsubsection{Answer Relevancy}

答案是否真正回答了问题，而不是泛泛而谈。


\subsubsection{Context Precision}

提供给模型的上下文中，有多少内容是真正相关的。这个指标能帮助你发现``检索很多，但噪声更大''的问题。


\subsection{7.7.3 RAGAS}

RAGAS 是一个很常见的 RAG 评估框架，可以对 retrieval 和 generation 做自动化评测。它不是万能的，但非常适合建立第一版离线评估基线。


工程建议：


\begin{itemize}[leftmargin=2em]
\item 建立 100-300 条高质量问答评测集
\item 每次调整 chunk、embedding、retriever、prompt 后都跑一遍
\item 把 Recall@K、Faithfulness、人工胜率一起看
\end{itemize}

\bigskip\hrule\bigskip

\section{7.8 常见失败模式与解决方案}

\subsection{7.8.1 检索到了不相关文档}

典型表现：模型回答有理有据，但完全答偏。


排查顺序：


\begin{enumerate}[leftmargin=2em]
\item chunk 是否过大或过小
\item embedding 模型是否不适合该语言/领域
\item 是否缺少 metadata filter
\item TopK 是否太大导致噪声多
\end{enumerate}

\subsection{7.8.2 检索到了正确文档，但答案仍然错}

说明错误更可能在生成阶段：


\begin{itemize}[leftmargin=2em]
\item prompt 没有要求``仅依据上下文''
\item 上下文排序不佳
\item 没做 rerank
\item chunk 太碎，模型无法拼出完整答案
\end{itemize}

解决思路：


\begin{itemize}[leftmargin=2em]
\item 加强 grounded prompt
\item 引入 reranker
\item 做 parent-child retrieval
\item 让模型给出引用片段
\end{itemize}

\subsection{7.8.3 明明有资料，但系统总说找不到}

这往往不是模型问题，而是覆盖率问题：


\begin{itemize}[leftmargin=2em]
\item 文档根本没入库
\item OCR 失败
\item 代码文件被排除
\item 查询改写不够好
\end{itemize}

解决：


\begin{itemize}[leftmargin=2em]
\item 加入多源数据
\item 做 ingestion 监控
\item 对空召回做 fallback（BM25、站内搜索、人工 FAQ）
\end{itemize}

\bigskip\hrule\bigskip

\section{7.9 一个完整的 Python 端到端 RAG 示例（LangChain）}

下面是一个最小但完整的 RAG 示例：加载 Markdown 文档、切块、embedding、存入 Chroma、检索并回答。


\subsection{7.9.1 安装依赖}

\begin{codeblock}
pip install langchain langchain-openai langchain-community chromadb tiktoken
\end{codeblock}

\subsection{7.9.2 准备文档}

假设 \texttt{./docs} 目录下有若干 \texttt{.md} 文件。


\subsection{7.9.3 完整代码}

\begin{codeblock}
from __future__ import annotations

from pathlib import Path

from langchain.text_splitter import RecursiveCharacterTextSplitter
from langchain_community.document_loaders import TextLoader
from langchain_community.vectorstores import Chroma
from langchain_openai import ChatOpenAI, OpenAIEmbeddings


DOC_DIR = Path("./docs")


def load_documents():
    docs = []
    for path in DOC_DIR.glob("*.md"):
        loader = TextLoader(str(path), encoding="utf-8")
        docs.extend(loader.load())
    return docs


def build_vectorstore():
    docs = load_documents()
    splitter = RecursiveCharacterTextSplitter(
        chunk_size=500,
        chunk_overlap=80,
        separators=["\n## ", "\n### ", "\n", "。", " ", ""],
    )
    chunks = splitter.split_documents(docs)

    vectorstore = Chroma.from_documents(
        documents=chunks,
        embedding=OpenAIEmbeddings(model="text-embedding-3-large"),
        persist_directory="./rag_chroma",
    )
    return vectorstore


def answer_question(question: str) -> str:
    vectorstore = build_vectorstore()
    retriever = vectorstore.as_retriever(search_kwargs={"k": 4})
    retrieved_docs = retriever.invoke(question)

    context = "\n\n".join(doc.page_content for doc in retrieved_docs)

    llm = ChatOpenAI(model="gpt-4o", temperature=0)
    prompt = f"""
你是一名企业知识库问答助手。
你只能依据给定上下文回答问题。
如果上下文中没有答案，请明确回答``根据当前知识库无法确定''。

上下文：
{context}

问题：
{question}
""".strip()

    resp = llm.invoke(prompt)
    return resp.content


if __name__ == "__main__":
    print(answer_question("系统为什么要使用 KV Cache？"))
\end{codeblock}

\subsection{7.9.4 这段代码缺了什么}

它能跑，但离生产还很远。你还需要：


\begin{itemize}[leftmargin=2em]
\item 文档增量更新
\item 检索缓存
\item rerank
\item 引用来源
\item 失败重试
\item 评估集
\item prompt injection 防护
\item 监控：空召回率、引用命中率、回答满意度
\end{itemize}

\bigskip\hrule\bigskip

\section{7.10 一个更接近生产的 RAG 设计建议}

如果你要设计企业级 RAG，建议按下面的结构拆分模块：


\begingroup
\scriptsize
\setlength{\tabcolsep}{3pt}
\begin{longtable}{|p{0.47\textwidth}|p{0.47\textwidth}|}
\hline
模块 & 建议职责 \\ \hline
Ingestion Service & 拉取 PDF/HTML/代码/数据库内容并清洗 \\ \hline
Chunking Service & 负责分块、标题路径、版本号记录 \\ \hline
Embedding Service & 统一 embedding 模型与缓存 \\ \hline
Retrieval Service & dense/sparse/hybrid 检索与 metadata filter \\ \hline
Rerank Service & 重排候选结果 \\ \hline
Prompt Builder & 组装 system prompt、context、citations \\ \hline
Answering Service & 调用 LLM 并输出结构化答案 \\ \hline
Evaluation Pipeline & 离线评测、A/B 实验、回归检查 \\ \hline
\end{longtable}
\endgroup

这套分层的好处在于：你以后换 embedding 模型、换向量库、换 reranker 时，不需要重写整个应用。


\bigskip\hrule\bigskip

\section{7.11 面试高频问法：如何回答更像工程师}

\subsection{问：为什么不能直接把所有文档塞进长上下文？}

答法要点：


\begin{itemize}[leftmargin=2em]
\item 成本高、延迟高
\item 上下文噪声会污染答案
\item 1M context 并不意味着 1M token 都同样可用
\item 检索本质是在做``信息压缩和定位''
\end{itemize}

\subsection{问：RAG 效果差时你先查什么？}

优秀回答：


\begin{enumerate}[leftmargin=2em]
\item 先区分是 retrieval 问题还是 generation 问题
\item 看评测集上的 Recall@K
\item 抽样查看 chunk 和召回结果
\item 再决定改 embedding、chunking、rerank 还是 prompt
\end{enumerate}

\subsection{问：Hybrid retrieval 为什么常优于纯向量检索？}

因为企业数据里大量关键信息是编号、函数名、报错码、表名、SKU 之类的精确 token，纯 dense 检索对这些信号未必敏感。


\bigskip\hrule\bigskip

\section{7.12 Query Transformation、引用与线上监控}

很多 RAG 新手会直接拿``用户原话''去检索，但真实业务里的问题经常包含口语、省略、缩写和上下文依赖。比如：


\begin{quote}
``昨天那个退款接口又慢了，会不会还是 MQ 堆积？''\\
\end{quote}

如果直接检索，这句话里真正有用的关键信号其实很散。于是我们需要 query transformation（查询改写）。


\subsection{7.12.1 常见改写方式}

\begin{enumerate}[leftmargin=2em]
\item Rewrite：把口语改成完整技术问题
\item Expand：补充同义词、英文术语、缩写全称
\item Decompose：拆成多个子问题
\item Normalize：抽取时间、服务名、错误码等结构化条件
\end{enumerate}

一个更可检索的改写结果可能是：


\begin{itemize}[leftmargin=2em]
\item 退款接口超时原因
\item MQ backlog 导致 refund API latency
\item payment-service refund timeout queue congestion
\end{itemize}

成熟系统通常不会只用改写后的 query，而是``原 query + 改写 query 并行检索，再合并去重''。这样既保留原始关键词，又提高语义覆盖率。


\subsection{7.12.2 为什么引用很重要}

企业用户不只想知道答案，还想知道``依据是什么''。因此好的 RAG 回答通常包含：


\begin{itemize}[leftmargin=2em]
\item 最终答案
\item 引用来源
\item 关键片段
\item 不确定性说明
\end{itemize}

例如：


\begin{codeblock}
答案：支付超时更可能与队列堆积有关。
依据：
1. 《退款服务故障复盘》2026-05-18
2. payment-service dashboard 周报第 3 节
不确定性：若你们已切换到 v2 队列配置，需要进一步核对最新部署单。
\end{codeblock}

这不仅降低幻觉风险，也方便用户人工复核。


\subsection{7.12.3 RAG 上线后的核心监控指标}

一套最低可用的 RAG 监控面板，建议至少有：


\begingroup
\scriptsize
\setlength{\tabcolsep}{3pt}
\begin{longtable}{|p{0.47\textwidth}|p{0.47\textwidth}|}
\hline
指标 & 说明 \\ \hline
empty\_retrieval\_rate & 空召回比例 \\ \hline
avg\_context\_tokens & 平均送入模型的上下文长度 \\ \hline
citation\_hit\_rate & 带有效引用的回答比例 \\ \hline
fallback\_rate & ``无法确定''或转人工比例 \\ \hline
index\_freshness\_lag & 文档更新到索引生效的延迟 \\ \hline
\end{longtable}
\endgroup

离线 Recall@K 很高，不代表线上体验就好。因为线上还会受到多轮上下文、文档新鲜度、缓存命中和用户提问风格的影响。


\bigskip\hrule\bigskip

\section{7.13 Agentic RAG 与多跳检索}

简单 RAG 通常只检索一次，但复杂问题常常需要多跳（multi-hop）推理。例如：


\begin{quote}
``为什么上海区用户昨天支付失败率上升，而退款成功率没有同步下降？''\\
\end{quote}

这个问题可能需要依次检索：


\begin{enumerate}[leftmargin=2em]
\item 支付服务告警记录
\item 上海区流量变更公告
\item 退款服务错误率报表
\item 最近一次配置变更单
\end{enumerate}

Agentic RAG 的含义是：让 Agent 决定是否继续检索、检索哪些子问题、何时停止。它比固定单次检索更强，但也更贵、更难控。所以通常需要：


\begin{itemize}[leftmargin=2em]
\item 工具调用次数上限
\item 检索回路检测
\item 每轮检索后的摘要压缩
\item 成本与延迟预算
\end{itemize}

如果你在面试里能把`\texttt{单次检索 RAG''和}`多跳 Agentic RAG''区别清楚，基本就超过了只会搭 demo 的候选人。


\bigskip\hrule\bigskip

\section{7.14 RAG 安全、权限与数据新鲜度}

很多团队以为 RAG 的主要风险是``答错''，实际上企业环境里还有三个更现实的问题：越权、脏数据和过期数据。


\subsection{7.14.1 检索越权}

如果你的知识库面向多个角色，例如普通员工、主管、管理员，那么``是否能检索到某段内容''本身就是权限问题。不要指望模型在回答时再自己判断哪些内容不能说，正确做法是：


\begin{itemize}[leftmargin=2em]
\item 在检索前做 ACL / role filter
\item metadata 中明确 \texttt{access\_level}
\item 高敏感内容默认不入通用索引
\end{itemize}

\subsection{7.14.2 脏数据污染}

如果源文档本身过时、矛盾或内容复制粘贴错误，RAG 会把这些问题忠实放大。因此 ingestion 阶段最好有：


\begin{itemize}[leftmargin=2em]
\item 文档版本字段
\item 生效时间和失效时间
\item 来源可信度等级
\item 重复文档检测
\end{itemize}

\subsection{7.14.3 数据新鲜度}

很多产品投诉``AI 总回答旧流程''，本质不是模型落后，而是索引刷新慢。一个成熟系统会对以下延迟做监控：


\begin{enumerate}[leftmargin=2em]
\item 文档变更时间到抓取时间
\item 抓取时间到切块完成时间
\item 切块完成时间到索引可查询时间
\end{enumerate}

如果这条链路总延迟是 6 小时，那么你就不能对业务方承诺``分钟级更新''。


\bigskip\hrule\bigskip

\section{7.15 一套实用的 RAG 调优顺序}

RAG 系统效果不好时，很多团队第一反应是换模型。其实更高效的调优顺序通常是：


\begin{enumerate}[leftmargin=2em]
\item 先检查文档是否真的入库了
\item 再检查 chunk 是否保留了结构边界
\item 再看 Recall@K 是否足够
\item 然后引入 hybrid retrieval
\item 再尝试 rerank
\item 最后再调 prompt 和生成模型
\end{enumerate}

原因很简单：如果召回阶段就把正确证据漏掉了，后面的模型再强也救不回来。


\subsection{7.15.1 一个典型案例}

某内部技术知识库最初方案：


\begin{itemize}[leftmargin=2em]
\item chunk size 1000
\item 纯 dense 检索
\item 无 rerank
\item 直接将 top 8 chunk 塞给模型
\end{itemize}

问题表现：


\begin{itemize}[leftmargin=2em]
\item Recall@5 只有 0.63
\item 回答常引用不相关段落
\item 平均上下文 4200 tokens，成本高
\end{itemize}

优化后方案：


\begin{itemize}[leftmargin=2em]
\item chunk size 改为 350 + 50 overlap
\item 增加 BM25 混合检索
\item 先召回 20 条，再 rerank 取前 5
\item 回答时强制引用来源
\end{itemize}

结果往往会变成：


\begin{itemize}[leftmargin=2em]
\item Recall@5 提升到 0.81
\item Faithfulness 明显改善
\item 平均上下文降到 1900 tokens
\end{itemize}

这类案例最能体现 RAG 是系统工程，而不是单点技巧。


\subsection{7.15.2 什么时候不该用 RAG}

不是所有 AI 问答都要上 RAG。以下场景就不一定合适：


\begin{enumerate}[leftmargin=2em]
\item 纯结构化查询，数据库直接答更准
\item 实时性极强且数据秒级变化，用 API/工具调用比索引检索更靠谱
\item 问题高度固定，FAQ 模板和规则引擎已经足够
\end{enumerate}

RAG 的优势是把`\texttt{非结构化知识''接入生成系统，而不是替代所有数据访问方式。很多成熟 Agent 实际采用的是}`RAG + Tool Use''的混合架构：知识类问题走 RAG，事务类问题走 API，分析类问题再由模型汇总。


\subsection{7.15.3 一个上线前检查清单}

在真正把 RAG 发布给用户之前，建议至少确认：


\begin{enumerate}[leftmargin=2em]
\item 有一套可复现的离线评测集
\item 空召回时有明确 fallback 策略
\item 回答能附带引用来源
\item 高权限文档不会被普通用户检索到
\item 文档更新到索引生效的延迟可监控
\end{enumerate}

RAG 最怕的不是``效果一般''，而是你根本不知道它什么时候开始失效。能监控、能回放、能回滚，才叫可运营的 RAG。


\subsection{7.15.4 让产品、算法和后端说同一种语言}

RAG 项目很容易失败在协作上：产品说`\texttt{回答不准''，算法说}\texttt{召回没问题''，后端说}`接口都通了''。为了避免这种各说各话，建议团队统一使用以下三层诊断语言：


\begin{enumerate}[leftmargin=2em]
\item 检索层：有没有把正确证据找回来
\item 排序层：找回来的证据是否排在前面
\item 生成层：模型是否忠于证据作答
\end{enumerate}

一旦团队能按这三层拆问题，RAG 的迭代速度会明显提升。

进一步说，RAG 项目的成功标准也应该分层定义：如果目标是降低人工客服转接率，就不能只看 Recall@K；还要看最终回答是否真的解决了用户问题，以及引用是否让用户敢于相信答案。

这也是为什么成熟团队会同时维护检索指标、生成指标和业务指标，而不是把所有责任都推给模型排行榜。

对转行面试者来说，如果你能把 RAG 讲成`\texttt{数据处理 + 检索排序 + 生成约束 + 评估运营''的完整闭环，而不是只会说}`向量数据库 + 大模型''，基本就已经具备了明显的竞争优势。

更进一步，如果你还能说出空召回、错误引用、版本过期、权限串漏这些真实故障模式，面试官通常会直接判断你做过真正的 RAG 项目。

而这，正是 AI Agent 岗位和普通``会调接口''岗位之间最关键的差别。

真正难的不是把 demo 跑通，而是让它在真实用户、真实数据和真实组织协作中持续稳定工作。

这才是工程化的核心。

也是企业愿意为你付薪水的原因。

而不是只奖励会写 demo 的人。

长期主义更重要。

确实。


\bigskip\hrule\bigskip

\section{本章要点}

\begin{enumerate}[leftmargin=2em]
\item RAG 的核心价值是补充新知识、接入私有数据、降低幻觉。
\item Naive RAG 只是起点，生产系统通常需要 query rewrite、hybrid retrieval、rerank 和评估。
\item 文档加载与 chunking 对效果影响极大，很多问题在进入模型前就已经埋下。
\item Dense 检索擅长语义，Sparse 检索擅长精确术语，Hybrid 往往是企业系统的默认选择。
\item Reranker 是把`\texttt{能召回''变成}`能排对''的关键步骤。
\item RAG 必须分层评估：Recall@K、MRR、nDCG、Faithfulness、Context Precision 缺一不可。
\item 诊断 RAG 系统时，首先要区分问题发生在检索阶段还是生成阶段。
\item 真正可维护的 RAG 架构应模块化，让 ingestion、retrieval、rerank、generation、evaluation 各自可替换。
\end{enumerate}

\section{延伸阅读}

\begin{enumerate}[leftmargin=2em]
\item LangChain、LlamaIndex 官方文档：重点看 retriever、reranker、evaluation 相关模块。
\item RAGAS 项目文档：理解自动评估指标的适用边界。
\item BM25、HNSW、Cross-Encoder reranker 相关资料：帮助你把``检索''从黑盒变成可调系统。
\item 各大向量数据库的 hybrid search 与 metadata filter 文档：这是企业落地中最容易产生差异的部分。
\item 实战练习：拿一套自己的项目文档，分别做 naive RAG 和 hybrid + rerank 版本，对比 Recall@K、回答准确率和 token 成本。
\end{enumerate}
